\subsection{Restriction of differentiable mappings}

Let n and r be positive integers, let U be an open subset of \(R^{n}\), and let F be a Banach space. Denote by \(\mathcal{C}^{r}(\mathrm{U};\mathrm{F})\) the vector space of maps of class \(C^{r}\) from U into F, endowed with the topology of \(C^{r}\) compact convergence. Recall that this is the supremum of the topologies of \(C^{r}\)-uniform convergence on K (VAR, R2, 12.3.10, p. 56), as K ranges over the set of compact subsets of U. The space \(\mathcal{C}^{0}(\mathrm{U};\mathrm{F})\) is none other than the space \(\mathcal{C}(\mathrm{U};\mathrm{F})\) of continuous maps from U into F, equipped with the topology of compact convergence.

Lemma 3. --- Let A be the set of multi-indices \(\alpha \in \mathbf{N}^n\) such that \(|\alpha| \leqslant r\) and let u be the linear map \(f \mapsto (\partial^\alpha f)_{\alpha \in \mathrm{A}}\) from \(\mathcal{C}^r(\mathrm{U};\mathrm{F})\) into \(\mathcal{C}(\mathrm{U};\mathrm{F})^{\mathrm{A}}\) .

\begin{enumerate}
\def\labelenumi{\alph{enumi})}
\item
  The map u is injective, continuous, strict, and its image is closed.
\item
  The topological vector space \(\mathcal{C}^{r}(\mathrm{U};\mathrm{F})\) is complete.
\end{enumerate}

The map \(u\) is linear and injective. It is continuous and strict by definition of the topology of \(\mathcal{C}^r (\mathrm{U};\mathrm{F})\) .

Let \((g_{\alpha})_{\alpha\in\mathrm{A}}\) be a point of \(\mathcal{C}(\mathrm{U};\mathrm{F})^{\mathrm{A}}\) adherent to the image of u. There exists a filter F on \(\mathcal{C}^{r}(\mathrm{U};\mathrm{F})\) such that one has \(g_{\alpha}=\lim_{f,\mathcal{F}}\partial^{\alpha}f\) in \(\mathcal{C}(\mathrm{U};\mathrm{F})\) for all \(\alpha\in A\). Let m be an integer such that \(0\leqslant m\leqslant r\) . By induction on m, it follows from th. 1 of FVR, II, p. 2, that the map \(g_{0}\) is of class \(C^{m}\) and that \(g_{\alpha}=\partial^{\alpha}g_{0}\) for all \(\alpha\in\mathbf{N}^{n}\) with \(|\alpha|\leqslant m\) . Thus \(g_{0}\) belongs to the space \(\mathcal{C}^{r}(\mathrm{U};\mathrm{F})\) and its image under u is \((g_{\alpha})_{\alpha\in\mathrm{A}}\) . This proves that the image of u is closed, and thus establishes assertion a).

The assertion \(a\) ) implies that the space \(\mathcal{C}^{r}(\mathrm{U};\mathrm{F})\) is isomorphic to its image in \(\mathcal{C}(\mathrm{U};\mathrm{F})^{\mathrm{A}}\) ; since this image is closed and the space \(\mathcal{C}(\mathrm{U};\mathrm{F})^{\mathrm{A}}\) is complete (TG, X, p. 9, cor. 3 of th. 2 and TG, II, p. 17, prop. 10), the space \(\mathcal{C}^{r}(\mathrm{U};\mathrm{F})\) is complete.

PROPOSITION 8. --- Suppose that F is finite-dimensional. Let s be an integer such that \(0 \leqslant s < r\) and let V be a relatively compact open subset of U. The linear map \(f \mapsto f|V\) from \(\mathcal{C}^{r}(U;F)\) into \(\mathcal{C}^{s}(V;F)\) is compact.

Let W be the set of functions \(f \in \mathcal{C}^{r}(\mathrm{U};\mathrm{F})\) whose partial derivatives of order \(\leqslant r\) take at every point of \(\overline{V}\) a value of norm less than 1. The set W is a neighborhood of 0 in the space \(\mathcal{C}^{r}(\mathrm{U};\mathrm{F})\) . Let \(\alpha\) be a multi-index such that \(|\alpha| \leqslant s\) . Consider the set H of functions of the form \((\partial^{\alpha}f)|\mathrm{V}\) for f in W. The set H is an equicontinuous subset of \(\mathcal{C}(\mathrm{V};\mathrm{F})\) by the mean value theorem (VAR, R, 2.2.3). Moreover, for every \(x \in V\) , the image of H under the map \(g \mapsto g(x)\) is a bounded subset, hence relatively compact, of F. By the Ascoli theorem (TG, X, p. 18, cor. 2), the set H is relatively compact in \(\mathcal{C}(\mathrm{V};\mathrm{F})\) . This proves that the linear map \(f \mapsto (\partial^{\alpha}f)|\mathrm{V}\) from \(\mathcal{C}^{r}(\mathrm{U};\mathrm{F})\) into \(\mathcal{C}(\mathrm{V};\mathrm{F})\) is compact. The proposition then follows from Lemma 3 and Remarks 5 and 6 of III, p. 3.

